
noting that @@M0@@ is a \emph{perfect} morphism, since it is obtained by the flat base change @@M1@@ from the morphism @@M2@@, which is perfect because @@M3@@ is regular. One thus obtains a quotient group
\sgaStableTarget{sga6.e10.eq.7-11-3}@@M4@@
of @@M5@@. This group is moreover naturally a module over @@M6@@, since one sees at once from the definitions that the product of an element of @@M7@@ with an element of @@M8@@ is algebraically equivalent to zero when one of the two factors is, using (IV 2.12.2) for @@M9@@.

One defines (when working with naive @@M10@@) the ideal of @@M11@@ of elements algebraically equivalent to zero as the ideal generated by the @@M12@@, where @@M13@@ and @@M14@@ is algebraically equivalent to zero; the quotient graded ring
\sgaStableTarget{sga6.e10.eq.7-11-4}@@M15@@
is then precisely the Chern ring associated with the @@M16@@-ring @@M17@@.

Finally, one defines the graded subgroup of @@M18@@ of elements algebraically equivalent to zero (@@M19@@ being assumed locally noetherian) as the subgroup whose component of dimension @@M20@@ is the image in @@M21@@ of the subgroup of @@M22@@ consisting of the elements algebraically equivalent to zero in @@M23@@. One finds at once that the quotient graded group
\sgaStableTarget{sga6.e10.eq.7-11-5}@@M24@@
is a module over the ring @@M25@@, by passage to the quotient from the module structure of @@M26@@ over @@M27@@ (3.1).

\sgaStableTarget{sga6.e10.para.7-11-6}\textbf{7.11.6.} The preceding definitions would still make sense without assuming that @@M28@@ is algebraically closed, taking @@M29@@ above to be \emph{geometrically} connected over @@M30@@. It is not clear, however, that these definitions would be very useful; in particular, it does not appear that, in the case of the group @@M31@@, one
